% chapters/ch17-requirements-and-scope.tex
\chapter{Requirements and Scope}

\begin{learningobjectives}
In this chapter you will learn how to define scope for a conceptual design.
You will be able to separate requirements that shape meaning from those that belong to later implementation.
You will practice writing clear scope boundaries and assumptions.
\end{learningobjectives}

\section{Why scope matters}
A model is only as clear as its boundaries.
If the scope is too wide, the model becomes vague.
If it is too narrow, it becomes brittle.
Scope is not a technical detail.
It is a design decision about meaning.

\section{What belongs in conceptual design}
Conceptual design captures facts, rules, and relationships.
It answers: what exists, what it means, and what is allowed.
It does not answer how data will be stored or optimized.

\subsection{Entities, attributes, and relationships}
Entities represent identity.
Attributes describe those entities.
Relationships express associations and constraints.
This is the core vocabulary of the model.

\subsection{Rules and constraints}
Rules belong in the conceptual model when they define valid states.
A model without constraints is a story, not a design.

\begin{examplebox}
\textbf{In-scope statements}\
A member is uniquely identified by \texttt{member\_id}.\\
A loan must reference exactly one member and one copy.\\
A copy may have at most one active loan at a time.
\end{examplebox}

\section{What does not belong yet}
Some topics are important, but not conceptual.
They are deliberately out of scope in this phase.

\subsection{Performance and storage}
Indexes, caching, and partitioning are implementation concerns.
They depend on workload and technology choices.
They do not change meaning.

\subsection{SQL and physical schemas}
SQL expresses the model, but it is not the model.
The conceptual design should be complete without a specific dialect.

\subsection{User interfaces and workflows}
UI workflows influence data needs, but they are not rules of the domain.
Keep the model neutral and stable.

\begin{examplebox}
\textbf{Out-of-scope statements}\
“Add an index on member\_id.”\\
“Use a JSON column for tags.”\\
“Show a drop-down for status.”
\end{examplebox}

\section{Sources of requirements}
Requirements come from people, documents, and existing systems.
No single source is enough.
You need triangulation.

\subsection{Stakeholder interviews}
Listen for nouns, verbs, and constraints.
Ask for exceptions and edge cases.
Those define the shape of the model.

\subsection{Policies and artifacts}
Policies describe rules.
Forms and reports reveal attribute expectations.
Existing databases reveal legacy assumptions.

\section{Writing requirements as rules}
A good rule is testable and specific.
It avoids “usually” and “typically”.
It uses explicit cardinality and optionality.

\begin{principlebox}
\textbf{Principle}\
Model meaning, not implementation.
\end{principlebox}

\section{Defining boundaries and assumptions}
State what you are not modelling.
State what you assume to be true.
This makes the model stable and reviewable.

\subsection{Boundary statements}
Examples include regions, time ranges, and product lines.
They reduce ambiguity and prevent scope creep.

\subsection{Assumption log}
Assumptions should be explicit and revisited.
They are part of the model’s contract.

\begin{chaptersummary}
Scope is a design decision, not a formality.
Conceptual design includes entities, attributes, relationships, and rules.
It excludes SQL, performance tuning, and UI decisions.
Clear boundaries make the model coherent and reviewable.
\end{chaptersummary}

\begin{exercisebox}
Define scope for a small domain of your choice.
\begin{enumerate}[label=\arabic*.]
  \item Write five in-scope rules as precise statements.
  \item List three out-of-scope concerns you will not model yet.
  \item Identify two sources you would use to validate the rules.
  \item Write three assumptions that might change later.
  \item Explain how you would communicate scope to stakeholders.
\end{enumerate}
\end{exercisebox}
